\section{Conclusion}

A broader technical standard can improve interoperability and diffuse common technological improvements across products. But it can also change what firms choose to innovate. This paper isolates that second effect by allowing competing firms to allocate a fixed R\&D capacity between common-layer and proprietary innovation after the regulator chooses the common scope.

The private and social responses move in opposite directions. Firms reduce common-layer R\&D as the standard expands because their improvements increasingly benefit a rival, whereas a social planner increases common-layer R\&D because the same transferability broadens the social return to innovation. This innovation-composition distortion feeds back into the regulator's policy problem. Complete standardization is optimal for every fixed positive common-R\&D allocation and under the first best, but sufficiently strong product-market rivalry yields a unique interior standardization scope when R\&D allocation is decentralized.

The policy interpretation is explicitly second-best. Selective standardization is not desirable because interoperability is intrinsically costly in the model. It is desirable because standardization is the regulator's only instrument and broader scope weakens private incentives to invest in the common layer. If another policy instrument could restore the private return to common innovation, the case for restricting the standard could disappear. Evaluating such instruments is beyond the frozen model.

The broader implication is that the scope of a technical standard and the direction of innovation should be analyzed jointly. A regulator that evaluates standardization while holding the innovation portfolio fixed may favor complete standardization even in an environment where the equilibrium response of firms makes a narrower common scope optimal. Future empirical work could test this mechanism by tracing R\&D or patent portfolios before and after changes in the scope of common technical interfaces, with particular attention to variation in downstream product substitutability.
